\documentclass{article}
\usepackage{graphicx} % Required for inserting images
\usepackage{amsmath}

\title{Ex 5.5}
\author{ERIK LAHUERTA ARTIEDA}
\date{November 2026}

\begin{document}

\maketitle

\subsection{Preliminaries \& Notation}
We will be drawing our uniform samples from a uniform distribution $U(0,1) \sim r_{1}, r_{2}$.\newline \newline
$w$ is our vector representation on the sphere, given by the convention $(r,\theta, \varphi)$, where r is the radius, $\theta$ is the polar angle, and $\varphi$ is the azimuthal angle. Since we are only interested in directions, r is not relevant. Then our differentials follow:
\[dw = sin(\theta)\space d\theta d\varphi\]
Similarly, because we are interested in $p(w)$, but want to convert to our angles, it also follows that:
\[p(w)=p(\theta, \varphi)\space sin(\theta)\] \newline
On cartesian coordinates, the conversion from polar to cartesian is as follows:
\[x=r\sin(\theta)\cos(\varphi)\]
\[y=r\sin(\theta)\sin(\varphi)\]
\[z=r\cos(\theta)
\]
These properties are used throughout the rest of the exercise and are useful to keep in mind.

\subsection{Uniformly Sampling a Hemisphere}
We will use inversion sampling. Sampling uniformly across the hemisphere means that: $p(w)=c$ for all $w$. Because this is a density function, it must mean that the integral over the domain must be equal to 1. Then:
\begin{align*}
\int_{H^{2}}p(w) dw
&=\int_{0}^{2\pi}\int_{0}^{2/\pi}c\ sin(\theta)\ d\theta d\varphi \\
&= c\int_{0}^{2\pi}\int_{0}^{2/\pi}sin(\theta) \ d\theta d\varphi \\
&= c\int_{0}^{2\pi}\left[-\cos(\theta)\right]^{2/\pi}_{0} d\varphi \\
&= c\int_{0}^{2\pi}1\ d\varphi \\
&=c*2\pi=1; \\ 
p(w) &=\frac{1}{2\pi}\ ;\\
p(\theta, \varphi) &= \frac{sin(\theta)}{2\pi} 
\end{align*}

Now that we have the joint distribution, we must find the marginal and conditional distributions. In this case, we will condition $\varphi\vert\theta$, although the reverse is valid as well.

\[
p_{\theta}(\theta)=\int^{2\pi}_{0}p(\theta, \varphi)\ d\varphi = \int^{2\pi}_{0}\frac{\sin(\theta)}{2\pi}\ d\varphi=\left[\varphi\frac{\sin(\theta)}{2\pi} \right]^{2\pi}_{0} = \sin(\theta)
\]
\[
p_{\varphi\vert\theta}(\varphi\vert\theta)=\frac{p(\theta,\varphi)}{p(\theta)}=\frac{\frac{\sin(\theta)}{2\pi}}{\sin(\theta)}=\frac{1}{2\pi}
\]

Now that we have found the marginal and conditional densities, we must find their CDF's in order to apply the inversion method:

\[
P_{\theta}(\theta)=\int_{0}^{\theta} \sin(\theta)\ d\theta= 1-\cos(\theta)
\]
\[
P_{\varphi\vert\theta}(\varphi\vert\theta) = \int^{\varphi}_{0}\frac{1}{2\pi}\ d\varphi = \frac{\varphi}{2\pi}
\]
Now we must apply the inversion method in order to convert uniform variables from 0 to 1 into uniformly distributed samples in polar coordinates.
\[
P_{\theta}(\theta)=1-\cos(\theta)=r_{1};\ \theta=\cos^{-1}(1-r_{1})=\cos^{-1}(r_{1})
\]
\[
P_{\varphi\vert\theta}(\varphi\vert\theta)=\frac{\varphi}{2\pi}=r_{2};\ \varphi =2\pi r_{2}
\]
We can safely convert $1-r_{1}$ to $r_{1}$, since both uniformly go from 0 to 1, but reversed.
We now have a conversion from uniform variables into our polar angles. The last step is to convert from polar coordinates to cartesian coordinates. Since we are only sampling the direction, a simplifying step is saying that the radius is 1.
There are useful trigonometric properties for the sake of simplification. One of them is: $\sin(\cos^{-1}(x))=\sqrt{1-x^{2}}$. Then it follows that: $\sin(\theta)=\sqrt{1-r_{1}^{2}}$
Another useful property is $\cos(\cos^{-1}(x))=x. $ It then follows that: $\cos(\theta)=r_{1}$
\begin{align*} 
x&=\sqrt{1-r_{1}^{2}}\ *\cos(2\pi r_{2}) \\
y&=\sqrt{1-r_{1}^{2}}\ *\sin(2\pi r_{2}) \\
z&=r_{1}
\end{align*}
This is the final result that we need to implement the sampling.
\subsection{Cosine-Weighted Hemisphere Sampling}
This time, we would like to sample from a PDF $p(w)\propto\cos(\theta)$. Following the process from last section:
\begin{align*}
\int_{H^{2}}p(w)\  dw
&=\int_{0}^{2\pi}\int_{0}^{2/\pi}c\ \cos(\theta)\sin(\theta)\ d\theta d\varphi \\
&= c\int_{0}^{2\pi}\left[-\frac{1}{2}\cos^{2}(\theta)\right]^{2/\pi}_{0} d\varphi \\
&= c\int_{0}^{2\pi}0.5\ d\varphi \\
&=c*\pi=1; \\ 
p(w) &=\frac{1}{\pi}\ ;\\
p(\theta, \varphi) &= \frac{\cos(\theta)\sin(\theta)}{\pi} 
\end{align*}
Instead of computing marginal and conditional densities, this time we shall use \textit{Malley's method} to generate the cosine-weighted points. This method consists of generating uniform samples on the unit disk, projecting them to the hemisphere above, and computing the direction from the center of the sphere. 
Using this method, we can use the sampling that was used on the disk sampling: $(r,\theta)=(\sqrt{r_{1}}, 2\pi r_{2})$. This method is equivalent to the points that we would obtain from doing the inversion method, so the final samples end up being the same.
\end{document}
